\documentclass[10pt,a4paper]{amsart}
\usepackage[margin=19mm]{geometry}
\usepackage{amsmath,amssymb,mathtools}
\usepackage[scr=rsfs,cal=euler]{mathalfa}
\usepackage{xfrac}
\usepackage[unicode,hidelinks,bookmarks=false]{hyperref}
\newcommand{\ie}{i.e.\ }
\newcommand{\cf}{cf.\ }
\newcommand{\resp}{respectively\ }
\newcommand{\Subsection}{\subsection}
\providecommand{\nobreakdash}{\nobreak\hbox{-}}
\setlength{\emergencystretch}{2em}
\multlinegap0pt
\usepackage{bm}
\usepackage{bbm}
\usepackage{dsfont}
\usepackage{xspace}
\usepackage{cancel}
\usepackage{mathtools}
\usepackage{amsthm}
\makeatletter
\@ifundefined{thm}{\newtheorem{thm}{Thm}}{}
\@ifundefined{lemma}{\newtheorem{lemma}[thm]{Lemma}}{}
\makeatother
\newcommand{\Fcal}{\mathcal{F}}
\newcommand{\Z}{\mathbb{Z}}
\newcommand{\rH}{\mathrm{H}}
\title[On refined nonvanishing conjectures by Kurihara and Kolyvagi]{On refined nonvanishing conjectures by Kurihara and Kolyvagin}
\author{Francesc Castella, Takamichi Sano}
\date{Source: arXiv:2601.14504v1 (CC BY 4.0)}
\begin{document}
\maketitle

\noindent\emph{Source and licence.} On refined nonvanishing conjectures by Kurihara and Kolyvagin. Version arXiv:2601.14504v1 (CC BY 4.0), distributed under the Creative Commons Attribution 4.0 International licence, as recorded in the arXiv licence metadata for this exact version and shown on the abstract page https://arxiv.org/abs/2601.14504v1. Full rights notice: licenses/arxiv-2601.14504-CC-BY-4.0.txt.\par\smallskip

\noindent\textit{Editorial scope.} Counted unit: the Lemma whose source label is \texttt{lem:ord-BK} in the article (source0.tex lines 1467--1474, source label \texttt{lem:ord-BK}), delivered together with the article's own complete original proof of it (source0.tex lines 1476--1518, one \texttt{proof} environment of 43 lines). No mathematics is written, repaired or re-derived: statement and proof are the article's own text line for line. Required context retained uncounted: eq:def-Falpha (lines 1165--1171). Every definition, display and label of those lines is retained unchanged. Omitted from this package and not redistributed: 1--1164, 1172--1466, 1475--1475, 1519--1588. Nothing inside a retained excerpt is omitted or altered: every retained body is delivered exactly as the article's lines are written, so no omission receipt of any kind accompanies this package. Citations: the retained excerpts carry no citation command at all, so no bibliography entry of the article is transcribed and no reference is redistributed. Reference adaptation: none was needed and none was made; every reference inside the delivered unit resolves inside it, so no unresolved reference or citation is produced. Printed slips kept literal: no printed word, symbol, hypothesis or number of the counted statement or of its proof was altered. Layout and class: the article's own class is \texttt{amsart} with packages including inputenc, amsmath, amsfonts, amssymb, graphicx, geometry, amsthm, lipsum, stmaryrd, dsfont, comment, xy, tikz-cd, hyperref; the wrapper is the lane's pinned \texttt{amsart} editorial wrapper at 10pt on A4, which restates the theorem-like environments the retained text needs and adds the article's own macro definitions that the retained text needs (\texttt{\textbackslash Fcal}, \texttt{\textbackslash Z}, \texttt{\textbackslash rH}), with extra wrapper packages (bm, bbm, dsfont, xspace, cancel, mathtools). The delivered numbering is the unit's own. Assets: no retained span contains a figure environment, a graphics-inclusion command, a PDF-page inclusion, a table or a TikZ picture; no image file is copied into this package, the package declares no asset dependency and no image file is redistributed with it. Licence: version arXiv:2601.14504v1 of this article is distributed under the Creative Commons Attribution 4.0 International licence, as recorded in the arXiv licence metadata for this exact version and shown on the abstract page https://arxiv.org/abs/2601.14504v1. A full rights notice is delivered at licenses/arxiv-2601.14504-CC-BY-4.0.txt. Characters: every retained excerpt is pure ASCII, so the delivered engine, XeLaTeX with its default Latin Modern text font, reports no missing character.
\begin{equation}\label{eq:def-Falpha}
\rH^1_{\Fcal_\alpha}(K_v,V(\alpha))=
\begin{cases}
\rH^1_{\rm ur}(K_v,V(\alpha)):=\rH^1(K_v^{\rm ur}/K_v,\rH^0(K_v^{\rm ur},V(\alpha)))&\textrm{if $v\nmid p$,}\\[0.2em]
{\rm im}\{\rH^1(K_v,V_v^+(\alpha))\rightarrow\rH^1(K_v,V(\alpha))\}&\textrm{if $v\mid p$,}
\end{cases}
\end{equation}

\begin{lemma}\label{lem:ord-BK}
Suppose $\alpha:\Gamma^{\rm ac}\rightarrow\Z_p^\times$  satisfies $\alpha\equiv 1\;({\rm mod}\,p^m)$ for some $m\gg 0$. Then for every prime $v$ of $K$ above $p$ we have
\[
\# \rH^0(K_v,A_v^-(\alpha^{\pm{1}}))\cdot\Z_p=\#\biggl(\frac{\rH^1_{\Fcal_\alpha}(K_v,T(\alpha^{\pm{1}}))}{\rH^1(K_v,T_v^+(\alpha^{\pm{1}}))}\biggr)\cdot\Z_p=
\#\widetilde{E}(\mathbb{F}_v)\cdot\Z_p
\]
for $m\gg 0$.
\end{lemma}

\begin{proof}
Let $v$ be a prime of $K$ above $p$. 
%In this proof we write $V(\alpha)^\pm={\rm Fil}_v^\pm(V(\alpha))$, $A(\alpha)^\pm={\rm Fil}_v^\pm(A(\alpha))$, etc. for the ease of notation. By \cite[Lem.~2]{flach-CT} % page 125
%we have the equality
%\begin{equation}\label{eq:flach-lem}
%\rH^1_f(K_v,V(\alpha))={\rm ker}\bigl\{\rH^1(K_v,V(\alpha))\rightarrow \rH^1(K_v^{\rm ur},V(\alpha)^-)\bigr\}.
%\end{equation}
%Since $p\nmid N$ by hypothesis, using inflation-restriction we ee that the map $\rH^1(K_v,V(\alpha)^-)\rightarrow \rH^1(K_v^{\rm ur},V(\alpha)^-)$ is injective. Hence from \eqref{eq:flach-lem} we deduce that $\rH^1_f(K_v,V(\alpha))\cong \rH^1(K_v,V(\alpha)^+)$, and so
Since $\rH^1_{\Fcal_\alpha}(K_v,A(\alpha))$ is the image of $\rH^1_{\Fcal_\alpha}(K_v,V(\alpha))$ in \eqref{eq:def-Falpha} under the natural map %on cohomology 
induced by $V(\alpha)\twoheadrightarrow A(\alpha)$, we see that
\begin{equation}\label{eq:f-div}
\rH^1_{\Fcal_\alpha}(K_v,A(\alpha))=\rH^1(K_v,A_v^+(\alpha))_{\rm div}.
\end{equation}
From the short exact sequence $0 \to T_v^+(\alpha) \to V_v^+(\alpha) \to A_v^+(\alpha) \to 0$ we obtain
\[
\frac{\rH^1(K_v,A_v^+(\alpha))}{\rH^1(K_v,A_v^+(\alpha))_{\rm div}}\cong \rH^2(K_v, T_v^+(\alpha)),
%\cong\rH^0(K_v,A(\alpha^{-1})^-)^\vee,
\]
which together with \eqref{eq:f-div} and local Tate duality amounts to
\[
\frac{\rH^1_{\Fcal_\alpha}(K_v,T(\alpha^{-1}))}{\rH^1(K_v,T_v^+(\alpha^{-1}))}\cong\rH^0(K_v,A_v^-(\alpha^{-1})),
\]
whence the first equality in the statement. The second equality is immediate from the fact that $A(\alpha^{\pm{1}})[p^m]\cong A[p^m]$ as $G_K$-modules, and so $\rH^0(K_v,A_v^-(\alpha^{\pm{1}}))\cong\rH^0(K_v,A_v^-)\cong \widetilde{E}(\mathbb{F}_v)[p^\infty]$.
%
%Now consider the commutative diagram with exact rows
%\[
%\xymatrix{
%0\ar[r]&\rH^0(K_v,A(\alpha)^+)\otimes\Q_p/\Z_p\ar[r]\ar@{->>}[d]&\rH^1(K_v,A(\alpha)^+)\ar[r]\ar[d]&\rH^1(K_v,A(\alpha)^+)[p^\infty]\ar[r]\ar[d]&0\\
%0\ar[r]&\rH^0(K_v,A(\alpha))\otimes\Q_p/\Z_p\ar[r]&\rH^1(K_v,A(\alpha))\ar[r] &\rH^1(K_v,A(\alpha))[p^\infty]\ar[r]&0,
%}
%\]
%where the surjectivity of the left vertical arrow follows  from the finiteness of $\rH^0(K_v,A(\alpha)^-)$. A simple diagram chase as in \cite[Prop.~2.5]{greenberg-cetraro} shows that the lower right horizontal map induces an isomorphism
%\begin{equation}\label{eq:chase}
%\frac{\rH^1(K_v,A(\alpha)^+)}{\rH^1(K_v,A(\alpha)^+)_{\rm div}}\cong \rH^1(K_v,A(\alpha)^-)[p^\infty]\cong \rH^0(K_v,A(\alpha)^-).
%\end{equation}
%(Note that this relies on the surjectivity of the natural map 
%\begin{equation}\label{eq:surj-red}
%\rH^0(K_v,A(\alpha))\rightarrow \rH^0(K_v,A(\alpha)^-).
%\end{equation}
%For $\alpha=1$, this amounts to the surjectivity for the reduction map $E(K_v)\rightarrow\widetilde{E}(\mathbb{F}_v)$; 
%(as used in \emph{loc.\,cit.}), 
%the surjectivity of \eqref{eq:surj-red} for $m\gg 0$ then follows immediately from that case using that $A(\alpha)[p^m]\cong A[p^m]$ as $G_K$-modules.) The result now follows from \eqref{eq:f-div}, and \eqref{eq:chase}, and local duality.
\end{proof}

\end{document}
